\section{Доказивање теорема у TLA+}

TLA+ не служи само за проверу модела \textit{(model checking)} тј. TLC проверава особине прегледом свих достижних стања,
што је изводљиво само за коначне моделе са конкретним вредностима константи.
За бесконачне системе или потпуно опште доказе користи се \textbf{TLAPS} (\textit{TLA+ Proof System}),
алат за формалне дедуктивне доказе директно у TLA+ синтакси \cite{tlapsProofTactics}.

TLAPS преводи доказе теорема на позадинске доказиваче теорема:
\begin{itemize}
    \item \textbf{Zenon} је аутоматски доказивач за логику првог реда,
    \item \textbf{Isabelle/TLA+} је интерактивни доказивач теорема,
    \item \textbf{SMT решавачи} нпр.\ Z3 за аритметичке леме.
\end{itemize}

\subsection{Инсталација TLAPS}
TLAPS долази као посебна инсталација и није део додатка за VS Code едитор. Потребно га је скинути
са званичног GitHub репозиторијума
\url{https://github.com/tlaplus/tlapm}. Када се скине инсталациони фајл, тада се .bin фајл за инсталацију
може локално инсталирати на следећи начин на десктопу
\begin{lstlisting}[language=bash]
./tlaps-*.bin -d $HOME/Desktop/tlaps
\end{lstlisting}
Када се заврши инсталација TLAPS можемо додати путању ка bin фолдеру како бисмо из тренутне сесије
терминала могли да покрећемо tlapm.
\begin{lstlisting}[language=bash]
export PATH=$PATH:~/Desktop/tlaps/bin
\end{lstlisting}
Да би проверили да ли TLAPS систем ради можемо покренути команду
\begin{lstlisting}[language=bash]
tlapm --version
\end{lstlisting}
Ако се појави број верзије система то значи да смо успешно инсталирали програм.
Напомињемо овде да је овде показана локална инсталација јер лако омогућава да се програм обрише без икаквих
системских измена. Ако се прекине тренутна сесија терминала морамо опет команду export, на тај начин
не дирамо системске путање. Овај приступ је одабран јер је релативно лаган да се уради без икакве шансе да се 
направи неки проблем. 

Могуће је повезати TLAPS са VS Code едитором и Toolbox-ом и покретати га тако или на начин који смо описали из терминала.

Када све ради како треба и када имамо нашу спецификацију са доказом неке теореме тада покренемо tlapm 
\begin{lstlisting}[language=bash]
tlapm <ime>.tla
\end{lstlisting}
излаз који добијамо у терминлу
\begin{lstlisting}[language=bash]
aleksandar@aleksandar-A520M-K-V2:~/Desktop/FormalneMetodeSkritpa/TLA$ tlapm LightSwitchProof.tla 
(* created new ".tlacache/TLAPS.tlaps/TLAPS.thy" *)
(* fingerprints written in ".tlacache/TLAPS.tlaps/fingerprints" *)
File "/home/aleksandar/Desktop/TLAPS/lib/tlaps/TLAPS.tla", line 1, character 1 to line 362, character 77:
[INFO]: All 0 obligation proved.
** Unexpanded symbols: ---

** Unexpanded symbols: STATE_TypeInvariant_, ACTION_Toggle_

** Unexpanded symbols: ---

** Unexpanded symbols: ---

(* created new ".tlacache/LightSwitchProof.tlaps/LightSwitchProof.thy" *)
(* fingerprints written in ".tlacache/LightSwitchProof.tlaps/fingerprints" *)
File "./LightSwitchProof.tla", line 1, character 1 to line 32, character 4:
[INFO]: All 12 obligations proved.
\end{lstlisting}
где нам на самом крају пише да су све обавезе доказане. Такође се генеришу и неки додатни фајлови са Isabelle доказом,
као што се може видети из порука са терминала.
\subsection{Структура доказа}

Доказ у TLA+ је \textbf{хијерархијски}. Разлаже се на нивое означене угластим заградама
\texttt{<1>}, \texttt{<2>}, \texttt{<3>} итд.\ Сваки корак се означава бројем, нпр.\ \texttt{<1>1}, \texttt{<1>2}.

Кључне конструкције:

\begin{center}
\begin{tabularx}{\textwidth}{|l|X|}
\hline
\textbf{Конструкција} & \textbf{Значење} \\
\hline
\texttt{THEOREM P} & исказ теореме коју треба доказати \\
\texttt{ASSUME H PROVE P} & доказ под претпоставком \texttt{H} \\
\texttt{BY facts} & закључи корак из наведених чињеница \\
\texttt{BY DEF Ime} & закључи разматрањем дефиниције \texttt{Ime} \\
\texttt{<1>n. QED} & завршни корак нивоа, обједињује претходне кораке у доказ\\
\texttt{PTL} & позива позадински доказивач за темпоралну логику \\
\hline
\end{tabularx}
\end{center}

\subsection{Пример: бројач са инваријантом}

Моделујемо бројач који почиње од нуле и увећава се за један у сваком кораку.
Желимо да формално докажемо да вредност бројача увек остаје природан број.
То је особина која важи за \emph{све} вредности бројача, без обзира на величину.

\begin{primer}[Доказ инваријанте бројача]

\textbf{TLA+ модул (\texttt{CounterProof.tla}):}
\begin{lstlisting}
---- MODULE CounterProof ----
EXTENDS Naturals, TLAPS

VARIABLE x

Init == x = 0
Next == x' = x + 1
Spec == Init /\ [][Next]_x

Inv == x \in Nat

THEOREM Spec => []Inv
<1>1. Init => Inv
      BY DEF Init, Inv
<1>2. Inv /\ [Next]_x => Inv'
      BY DEF Inv, Next
<1>3. QED
      BY <1>1, <1>2, PTL DEF Spec
====
\end{lstlisting}

Теорема каже: ако систем задовољава спецификацију \texttt{Spec}, онда \texttt{Inv} важи у
\emph{свим} стањима (темпорална формула \texttt{[]Inv}).

Доказ се одвија у три корака на нивоу \texttt{<1>}:
\begin{itemize}
    \item \texttt{<1>1} --- \textbf{базни случај}: ако је \texttt{Init} тачно, онда је \texttt{Inv} тачно.
    Разматрањем дефиниција: \texttt{Init} каже да је \texttt{x = 0}, а \texttt{Inv} каже да је \texttt{x \textbackslash in Nat} --- пошто је $0 \in \mathbb{N}$, следи тривијално.

    \item \texttt{<1>2} --- \textbf{индуктивни корак}: ако \texttt{Inv} важи у тренутном стању и
    систем изврши корак \texttt{[Next]\_x}, онда \texttt{Inv'} важи у следећем стању.
    Ако је \texttt{x \textbackslash in Nat} и \texttt{x' = x + 1}, онда је \texttt{x' \textbackslash in Nat} јер је скуп природних бројева затворен за сабирање.

    \item \texttt{<1>3. QED} --- обједињује \texttt{<1>1} и \texttt{<1>2} позивањем \texttt{PTL}
    (\textit{Propositional Temporal Logic}) доказивача, који из базног и индуктивног случаја аутоматски изводи \texttt{[]Inv}.
\end{itemize}

Ово је образац \textbf{индуктивне инваријанте}, најчешћа техника доказивања safety особина у TLA+:
\begin{enumerate}
    \item Докажи да инваријанта важи у почетном стању (\textit{база индукције}).
    \item Докажи да ако инваријанта важи пре прелаза, важи и после (\textit{индуктивни корак}).
    \item Из ова два корака TLAPS аутоматски изводи да инваријанта важи увек.
\end{enumerate}
Приметимо из овог примера да ово нисмо могли да докажемо провером модела. Променљива \texttt{x} је типа 
природних бројева и број стања је бесконачан и проверавач модела би се вртео заувек покушавајући да досегне 
и провери сва стања. У оваком слуачју морамо да посегнемо за напреднијим техникама и да формулишемо оваква
тврђења као теореме.
\end{primer}

\begin{primer}[Доказ безбедности прекидача за светло]
Враћамо се на пример прекидача за светло и формално доказујемо да \texttt{light} увек остаје
у скупу дозвољених вредности. Ово је особина коју смо раније само проверавали TLC-ом.

\textbf{TLA+ модул (\texttt{LightSwitchProof.tla}):}
\begin{lstlisting}
---- MODULE LightSwitchProof ----
EXTENDS TLAPS

VARIABLE light

Init == light = "off"

Toggle ==
    \/ /\ light = "off" /\ light' = "on"
    \/ /\ light = "on"  /\ light' = "off"

Next == Toggle

Spec == Init /\ [][Next]_light

TypeInvariant == light \in {"off", "on"}

THEOREM Spec => []TypeInvariant
<1>1. Init => TypeInvariant
      BY DEF Init, TypeInvariant
<1>2. TypeInvariant /\ [Next]_light => TypeInvariant'
  <2>1. ASSUME TypeInvariant, Toggle
        PROVE  TypeInvariant'
        BY <2>1 DEF TypeInvariant, Toggle
  <2>2. ASSUME TypeInvariant, UNCHANGED light
        PROVE  TypeInvariant'
        BY <2>2 DEF TypeInvariant
  <2>3. QED
        BY <2>1, <2>2 DEF Next
<1>3. QED
      BY <1>1, <1>2, PTL DEF Spec
====
\end{lstlisting}

Доказ следи исти образац индуктивне инваријанте, али са мање корака јер систем има само једну акцију.

\begin{itemize}
    \item \texttt{<1>1} у почетном стању \texttt{light = "off"}, па \texttt{light \textbackslash in \{"off", "on"\}} тривијално важи.
    \item \texttt{<1>2}  индуктивни корак се рашчлањује на два подслучаја:
    \begin{itemize}
        \item \texttt{<2>1}: акција \texttt{Toggle} мења \texttt{light} из \texttt{"off"} у \texttt{"on"} или обрнуто, у оба случаја вредност остаје у \texttt{\{"off", "on"\}},
        \item \texttt{<2>2}: ако \texttt{light} остаје непромењено, инваријанта тривијално важи.
    \end{itemize}
    \item \texttt{<1>3. QED}  \texttt{PTL} обједињује базни и индуктивни случај у \texttt{[]TypeInvariant}.
\end{itemize}
\end{primer}


Треба напоменути да провера модела (TLC) и доказивање теорема (TLAPS) се \textbf{допуњују}: TLC је бржи и аутоматски
проналази грешке у конкретним моделима, док TLAPS даје потпун доказ за све могуће вредности параметара.
Уобичајен ток корака је прво верификовати спецификацију са TLC-ом, а затим формално доказати кључне особине са TLAPS-ом.

